% ============================================================
%  4. ALCANCE   (guía, 2.4)
%  Propósito: definir hasta dónde llegará el proyecto y qué queda
%  FUERA. La guía advierte: evitar prometer resultados que no
%  puedan desarrollarse o evaluarse en el tiempo de la materia.
% ============================================================
\chapter{Alcance}
\label{cap:alcance}

El presente trabajo comprende la estimación del número esperado de consultas de ruta de
Trufi App por unidad territorial en el eje metropolitano de Cochabamba, desde la comprensión
y preparación de los datos hasta la selección, evaluación e interpretación de la técnica de
estimación y la formulación de una propuesta de uso para Trufi Association.

El estudio se delimita bajo las siguientes condiciones:

\begin{itemize}
  \item \textbf{Datos considerados.} Registro de consultas de Trufi App, estimaciones de
        población de Kontur Population y red de transporte mapeada en formato GTFS. No se
        emplean datos de viajes realizados, encuestas de movilidad ni información
        socioeconómica.
  \item \textbf{Período analizado.} Consultas registradas entre el 12 de septiembre de 2022
        y el 9 de junio de 2024.
  \item \textbf{Dominio territorial.} El eje metropolitano de Cochabamba y su entorno
        inmediato.
  \item \textbf{Unidad de análisis.} La celda hexagonal H3 de resolución 8, de
        aproximadamente 0,74~km\textsuperscript{2}.
  \item \textbf{Variable a estimar.} El número de consultas originadas en cada celda durante
        el período, en relación con su población residente, y la brecha entre lo observado
        y lo esperado.
  \item \textbf{Análisis incluidos.} Análisis exploratorio, construcción de la tabla de
        análisis por celda, comparación de técnicas de estimación con evaluación espacial,
        análisis de la brecha observado-esperado, contraste descriptivo con la red mapeada y
        análisis de sensibilidad.
  \item \textbf{Análisis excluidos.} La predicción de la evolución temporal de la demanda,
        la segmentación de usuarios, la inferencia de relaciones causales y la
        implementación operativa del resultado dentro de los sistemas de Trufi Association.
\end{itemize}

Los resultados están condicionados por limitaciones conocidas: las consultas reflejan el
comportamiento de quienes usan la aplicación y no el de toda la población; el registro no
permite saber si una consulta derivó en un viaje; Kontur Population es una estimación
modelada y no un censo; y la versión disponible del GTFS declara una vigencia posterior a
buena parte del período de las consultas. Por consiguiente, la solución constituye un
instrumento de apoyo a la priorización y no sustituye el criterio ni el conocimiento
territorial de los voluntarios.
